% ===========================================================================
\section{O que a representação sugere}\label{sec:perspectiva}
% ===========================================================================

\subsection{O dicionário em resumo}

O artigo construiu uma tradução e mediu suas consequências. A tabela abaixo
resume o que foi traduzido.

\begin{center}
\small
\begin{tabular}{ll}
\toprule
aritmética & reta (frequência) \\
\midrule
inteiro $m$ & ponto $\nfreq m=\log m$ \\
produto $mn$ & soma $\nfreq m+\nfreq n$ \\
$B$-suavidade & pertencer ao reticulado $\Lat(B)$ \\
ser composto & ser soma de geradores \\
ser primo & ser irredutível \\
divisibilidade por $p$ & pertencer a $\Lat(\{p\})$ \\
semente $\Seed_k$ & reticulado do bloco anterior \\
margem de classificação & $\rB_{\min}(k)=g(q_k)$ \\
\bottomrule
\end{tabular}
\end{center}

Nenhuma linha dessa tabela é profunda isoladamente: todas são o Teorema
Fundamental da Aritmética visto por logaritmos. O que a tabela oferece é uma
\emph{linguagem}, e a linguagem tem consequências que a notação multiplicativa não
tem: ela permite falar de distância, de resolução, de superposição e de contagem.
Os Teoremas~\ref{teo:exato} e \ref{teo:custo} e as
Proposições~\ref{prop:ganho} e \ref{prop:gemeo} são exemplos dessas consequências.

\subsection{Perguntas que ficam}

\paragraph{Outras bases.} A partição binária não é a única possível. Para uma base
$b\ge2$, o bloco $[b^k,b^{k+1})$ tem comprimento $b^k(b-1)$, e o menor fator de um
composto é $<b^{(k+1)/2}$; a semente correspondente é $\{p<b^k\}$. O análogo do
Teorema~\ref{teo:cota} vale com $b$ no lugar de $2$, mas a geometria muda: a
margem passa a depender de $b$, e não sabemos se o análogo do
Teorema~\ref{teo:exato} vale para $b>2$ — o vizinho $q+1$ deixa de estar sempre no
reticulado quando $q$ não é ímpar.

\paragraph{Outros crivos.} O crivo de Eratóstenes é o que a partição binária
organiza. O que acontece com crivos de roda (\emph{wheel sieves}), que riscam
múltiplos de vários primos de uma vez? A imagem correspondente seria uma soma de
tons em frequências de produtos, e a contagem coerente sugere que a altura dos
picos mediria quantos candidatos caem em cada classe residual.

\paragraph{Órbitas e a função de Dickman.} A contagem $k_r$ da
Proposição~\ref{prop:ganho} é o tamanho de uma órbita $\mathcal O_r$. Para $P$ fixo
e $m\le X$, essa contagem é o número de inteiros de $\mathcal O_r$ abaixo de $X$,
que é exatamente o objeto da teoria de números suaves e da função de
Dickman--de Bruijn $\Psi(X,y)$ \cite{tenenbaum}. A representação espectral dá a
$\Psi$ uma leitura: a altura do pico em $\freq r$ é $\Psi$-local.

\paragraph{A assinatura de pares gêmeos.} A Proposição~\ref{prop:gemeo} dá a lei
$q/T$ para o fundo. A constante ótima — a que aparece na
Tabela~\ref{tab:gemeos} como fator entre $q/T$ e o máximo medido, da ordem de
$1{,}8$ — não foi determinada. Determiná-la exigiria controlar a interferência
conjunta de todos os primos do sinal, e não só do vizinho imediato.

\paragraph{Aplicações visuais.} As figuras deste artigo foram feitas para serem
olhadas, mas cada uma delas é o resultado de um cálculo simples e curto. Um
instrumento interativo — em que o leitor escolhe $p$, o número de primos, a janela
e a faixa de leitura — permitiria percorrer a transição entre ``os pares se
fundem'' e ``os pares se separam'' que a Figura~\ref{fig:resolucao} mostra em dois
instantâneos. A Proposição~\ref{prop:janela} dá o critério exato para colorir essa
transição.

\subsection{O que este artigo não é}

Vale ser explícito sobre o alcance do que foi feito, para que o leitor saiba o que
esperar daqui.

A representação é fiel: os primos de um bloco aparecem como picos isolados, e
cada afirmação estrutural sobre eles tem uma tradução geométrica exata. Mas ela
não é um algoritmo competitivo. A aritmética decide em $O(\sqrt m)$ o que a
imagem exige uma janela de tamanho $\Theta(m)$ para mostrar
(Corolário~\ref{cor:local}). O valor do dicionário está em outro lugar: em tornar
visíveis relações que a notação multiplicativa esconde, e em dar a cada uma delas
uma lei quantitativa.

É, nesse sentido, o tipo de matemática que se faz para entender uma estrutura, e
não para calcular com ela. A estrutura em questão — a irredutibilidade dos
inteiros — é bem conhecida; o que é novo aqui é o ângulo, e o que o ângulo revela
são grandezas que não apareciam antes: a margem $g(q_k)$, o comprimento
$\tau(\Seed_k)$, o contraste $T/q$. Essas grandezas são as mesmas que governam
qualquer tentativa de ler aritmética por frequência, e é por isso que vale
registrá-las.
